\section{The observer-conditioned deal posterior}
\label{sec:belief}

\subsection{The hidden state is the initial deal}

In the dialect of \S\ref{sec:rules} a card changes hands in exactly one way, a
successful ask, which every player observes; declarations remove cards from
play, and that too is public. Hence:

\begin{proposition}[Public transfer]
\label{prop:transfer}
Assume the dialect of \S\ref{sec:rules}, in which the only card movements are
successful asks and declarations and both are announced. Let $c$ be a card whose
location has never been publicly revealed. Then $c$ is still held by the player to
whom it was dealt, the joint hidden state of the game at any time is exactly the
initial deal, and the current holder of every card is a deterministic function of
the initial deal and the public history.
\end{proposition}

\begin{proof}
A card changes hands only through a successful ask, and both the ask and its
outcome are announced, so every change of holder is a public event and a card
with no such event is where it was dealt. Replaying the public history from the
deal then determines current holdings, since each announced transfer moves a
named card between named seats and each declaration removes a named half-suit.
\end{proof}

Proposition~\ref{prop:transfer} eliminates the filtering recursion that
dominates inference in most imperfect-information card games: there is no
belief that decays and no dynamics to approximate, only a fixed hidden object
and an accumulating set of logical constraints on it. What follows is an
\emph{observer-conditioned deal posterior}: every player computes it from the
same public history but conditions additionally on their own private hand, so
the six posteriors differ and no single distribution is shared. It is exact
with respect to the rules-conditioned uniform prior of
\S\ref{sec:constraints}, and it is not a complete Bayesian opponent model,
since it conditions on nothing about how opponents choose among legal moves
(\S\ref{sec:policyprior}). Two interchangeable engines compute it, contrasted
in Figure~\ref{fig:architecture}: the exact reference engine of
\S\ref{sec:blocks}, and the approximate Fast path of \S\ref{sec:fast} that the
deployed \vfast{} policy runs.

\begin{figure}[ht]
\centering
\resizebox{\linewidth}{!}{% Pipeline of FishBot v0.4.  The dashed enclosure separates the exact reference
% engine (validation only) from the approximate path the deployed policy runs.
\begin{tikzpicture}[
  font=\small,
  node distance=0pt,
  box/.style   = {draw=fishgray!70, rounded corners=2pt, fill=white,
                  align=center, inner sep=5pt, minimum height=9mm, text width=30mm},
  inbox/.style = {box, fill=fishlight, draw=fishgreen!60},
  ref/.style   = {box, draw=fishgreen, dashed, fill=white},
  act/.style   = {box, fill=fishgreen!12, draw=fishgreen},
  lab/.style   = {font=\scriptsize\itshape, text=fishgray},
  ar/.style    = {-{Latex[length=2mm]}, draw=fishgray!85, thick},
  arref/.style = {-{Latex[length=2mm]}, draw=fishgreen, thick, dashed},
]

% --- inputs ---------------------------------------------------------------
\node[inbox] (hand)  at (0, 1.05) {observer's\\private hand};
\node[inbox] (pub)   at (0,-1.05) {public transcript\\asks, answers,\\declarations, counts};

% --- constraints ----------------------------------------------------------
\node[box, text width=34mm, minimum height=22mm] (con) at (4.2,0)
  {\textbf{rule constraints}\\[2pt]
   \scriptsize C1 own hand\\
   \scriptsize C2 revealed holders\\
   \scriptsize C3 exclusions\\
   \scriptsize C4 capacities\\
   \scriptsize C5 ask-legality certificates};

% --- the two inference paths ---------------------------------------------
\node[act, text width=36mm, minimum height=14mm] (fast) at (9.2, 1.55)
  {\textbf{Fast path} (deployed)\\[1pt]
   \scriptsize Sinkhorn capacity fitting\\
   \scriptsize + C5 conditioning\\
   \scriptsize $\hat\mu,\ \hat\tau,\ \hat\alpha$};

\node[ref, text width=36mm, minimum height=14mm] (block) at (9.2,-1.55)
  {\textbf{Reference engine}\\[1pt]
   \scriptsize half-suit block DP\\
   \scriptsize exact $Z,\ \mu,\ \tau,\ \alpha$\\
   \scriptsize validation and ablation only};

% --- decision -------------------------------------------------------------
\node[act, text width=30mm, minimum height=11mm] (ask) at (14.0, 1.55)
  {ask scorer\\\scriptsize \numAskFeats{} features + one-ply\\\scriptsize value, two-ply rerank};
\node[act, text width=30mm, minimum height=11mm] (decl) at (14.0,-0.15)
  {declaration rule\\\scriptsize pre-gates, then\\\scriptsize value-based stopping};
\node[act, text width=18mm, minimum height=9mm] (act) at (18.1, 0.70) {\textbf{action}};

% --- edges ----------------------------------------------------------------
\draw[ar] (hand.east)  -- ++(0.5,0) |- (con.west);
\draw[ar] (pub.east)   -- ++(0.5,0) |- (con.west);
\draw[ar] (con.east) -- ++(0.55,0) |- (fast.west);
\draw[arref] (con.east) -- ++(0.55,0) |- (block.west);
\draw[ar] (fast.east) -- (ask.west);
\draw[ar] (fast.south east) -- ++(0.45,0) |- (decl.west);
\draw[ar] (ask.east) -- ++(0.45,0) |- (act.west);
\draw[ar] (decl.east) -- ++(0.45,0) |- (act.west);
\draw[arref] (block.east) -- ++(0.7,0) |- ($(decl.south)+(0,-0.75)$)
      node[lab, below, pos=0.72] {substitution ablation (\S\ref{sec:ablations})} -- (decl.south);

% --- annotations ----------------------------------------------------------
\node[lab, align=center] at (9.2, 0.02) {both paths compute the\\same three quantities};
\node[lab, above] at (4.2, 1.35) {logical consequences of the rules};

\end{tikzpicture}
}
\caption{Inference and decision pipeline. The observer's private hand and the
public transcript determine the rule constraints C1--C5, from which two
interchangeable engines compute the same three quantities: per-card ownership
marginals, team-ownership probabilities, and named-allocation probabilities. The
\textbf{Fast path} (Sinkhorn capacity fitting with C5 conditioning) is what the
deployed \vfast{} policy runs, and produced every performance number reported in
this paper. The \textbf{reference engine} (exact half-suit block dynamic program)
is used for validation and for the substitution ablation only, and is not on the
deployed decision path.}
\label{fig:architecture}
\end{figure}

\subsection{Notation}
\label{sec:notation}

These symbols are used throughout the paper and are not varied. $\mathcal{H}$
is the set of live half-suits and $H \in \mathcal{H}$ one of them. Fix an
observer $i$. $U$ is the set of cards in live half-suits whose holder is not
publicly known to $i$, and for $c \in U$ the map $\sigma : U \to
\{0,\dots,5\}$ gives that card's holder, which by
Proposition~\ref{prop:transfer} is constant in time. $M_c \subseteq
\{0,\dots,5\}$ is the surviving support of $\sigma(c)$, and $q_p$ is the
capacity of seat $p$, the number of cards of $U$ that seat must hold. $Z$ is
the number of assignments $\sigma$ consistent with the constraints below, that
is the partition function of the posterior. The three quantities the policy
consumes are the per-card ownership marginal $\mu_{c,p} = \Pr[\sigma(c) = p]$;
the probability $\tau(H,T)$ that team $T$ owns every card of half-suit $H$;
and the probability $\alpha(A)$ that a named allocation $A$ --- an assignment
of all six cards of one half-suit to named seats --- is correct. Unhatted
symbols are values from the exact reference engine, hatted ones
($\hat\mu_{c,p}$, $\hat\tau(H,T)$, $\hat\alpha(A)$) the same values from the
Fast path of \S\ref{sec:fast}. Probabilities are conditional on the observer's
information $\mathcal{I}$ throughout; the conditioning bar is suppressed.

\subsection{The constraint system}
\label{sec:constraints}

The observer's information is exactly the following five constraints on
$\sigma$.

\begin{description}[style=unboxed,leftmargin=0pt,itemsep=4pt,parsep=0pt]
  \item[(C1) Own hand.] $\sigma(c) \ne i$ for every $c \in U$, and every card $i$
        holds is resolved.
  \item[(C2) Public transfers and correct declarations.] Any card whose holder
        has been revealed leaves $U$ with its holder known exactly. A successful
        ask reveals the holder \emph{before} the transfer, and that is the value
        relevant to every earlier constraint.
  \item[(C3) Exclusions.] An ask for $c$ proves the asker did not hold $c$; a
        miss proves the target did not hold $c$. Both are permanent, again by
        Proposition~\ref{prop:transfer}. The surviving support $M_c$ records
        them.
  \item[(C4) Capacities.] Hand counts are public, so for every player $p$
        \begin{equation}
          \bigl|\{c \in U : \sigma(c) = p\}\bigr| \;=\; q_p \;:=\; h_p - k_p ,
          \label{eq:capacity}
        \end{equation}
        where $h_p$ is $p$'s public card count and $k_p$ the number of live cards
        already known to be $p$'s. Note that $\sum_p q_p = |U|$ automatically,
        and that a declaration --- correct or not --- is absorbed by
        \eqref{eq:capacity} without special handling, because the count drop it
        causes is public.
  \item[(C5) Ask legality.] An ask for $c$ in half-suit $H$ by player $a$ proves
        that $a$ held some \emph{other} card of $H$ at that moment. Because
        unresolved cards never move, this is the time-independent disjunction
        \begin{equation}
          \bigvee_{d \in D} \bigl[\sigma(d) = a\bigr],
          \qquad D = \{d \in H \setminus \{c\} : d \text{ unresolved},\; a \in M_d\},
          \label{eq:disj}
        \end{equation}
        vacuous if some card of $H$ is already known to be $a$'s. The asking seat
        is written $a$ to keep $A$ reserved for the named allocation of
        \S\ref{sec:notation}. Every such certificate is confined to a single
        half-suit.
\end{description}

The prior is uniform over deals consistent with the rules, so the posterior is
uniform over the assignments $\sigma$ satisfying (C1)--(C5), and $\Pr[\sigma]
= 1/Z$ for each of them. This is exact with respect to that prior, and it is
policy-agnostic.

\subsection{Counting: a capacity-vector dynamic program}
\label{sec:counting}

Set (C5) aside for a moment. Counting assignments subject to (C3) and (C4) is a
degree-constrained bipartite counting problem --- the permanent of a $0/1$ matrix
with at most six distinct column types --- solved exactly by a forward--backward
recursion over the vector $n = (n_0,\dots,n_5)$ of \emph{remaining} capacities
from \eqref{eq:capacity}. Ordering $U$ as $c_1,\dots,c_{|U|}$, $F_j(n)$ counts
placements of the first $j$ cards leaving $n$ unused and $B_j(n)$ placements of
the rest out of $n$:
\begin{align}
  F_j(n) &= \sum_{p \,\in\, M_{c_j} :\, n_p < q_p} F_{j-1}(n + e_p),
  &F_0(q) &= 1, \label{eq:fwd}\\
  B_j(n) &= \sum_{p \,\in\, M_{c_{j+1}} :\, n_p \ge 1} B_{j+1}(n - e_p),
  &B_{|U|}(\mathbf{0}) &= 1, \label{eq:bwd}\\
  Z &= F_{|U|}(\mathbf{0}) \;=\; B_0(q), \label{eq:Z}
\end{align}
Here $e_p$ is the $p$-th unit vector. Recursions \eqref{eq:fwd}--\eqref{eq:Z}
count rather than sample, so they and the per-card marginal \eqref{eq:marg} of
\S\ref{app:layer}, which contracts the two tables at card $j$, are
\emph{exact} for (C1)--(C4); the layer index is fixed by the state, so a full
forward--backward pass costs at most $6 \times 10^{5}$ multiply--adds
regardless of the position (\S\ref{app:layer}, \S\ref{app:bound},
\eqref{eq:bound}). The backward table also yields a direct sampler,
rejection-free for (C1)--(C4) but combined with an exact accept/reject test
against \eqref{eq:disj} when a (C5) certificate is live, so the combined
procedure is exact but not rejection-free; it is the ground truth for the E2
cross-checks of \S\ref{sec:results-verify} (\S\ref{app:sampler}).

\subsection{Ask legality by half-suit block enumeration}
\label{sec:blocks}

Constraint (C5) is disjunctive and breaks the product form of
\eqref{eq:fwd}--\eqref{eq:bwd}, and inclusion--exclusion over the $k$
certificates a game accumulates --- tens of them --- would cost $2^k$ passes. The
escape is structural: every certificate lives inside one half-suit, so $U$
partitions into half-suit blocks admitting one of two exact treatments. A block with no live
certificate is expanded card by card as in \eqref{eq:fwd}--\eqref{eq:bwd}; one
with a live certificate is enumerated exhaustively --- at most
$\prod_{c \in H}|M_c| \le 5^6$ assignments --- each tested against
\eqref{eq:disj} directly, and the survivors aggregated into a count-vector
table
\begin{equation}
  g_H(t) \;=\; \bigl|\{\text{assignments } \beta \text{ of } H :
    \mathrm{count}(\beta) = t,\; \beta \models \text{certificates of } H\}\bigr|,
  \qquad \textstyle\sum_p t_p = m_H ,
  \label{eq:blocktable}
\end{equation}
where $m_H$ is the number of unresolved cards of $H$ and
$\mathrm{count}(\beta)_p$ is the number of them $\beta$ gives to seat $p$. Either
way a block consumes a known number of cards, so the layer identity survives and
the dynamic program runs over blocks rather than cards (\S\ref{app:groups} gives
the mechanics of both). Writing $S_t$ for the path mass of count vector $t$ at
its block's entry layer,
\begin{equation}
  S_t \;=\; \sum_{|n| \,=\, \text{entry layer}} F(n)\, B(n - t),
  \label{eq:pathmass}
\end{equation}
the three quantities of \S\ref{sec:notation} are read straight off the tables:
\begin{align}
  \mu_{c,p} &= \frac{1}{Z}\sum_t g_H^{(c \to p)}(t)\, S_t
  &&\text{(per-card ownership)} \label{eq:q1}\\
  \alpha(A) &= \frac{S_{t(A)}}{Z}
  &&\text{(the declaration query)} \label{eq:q2}\\
  \tau(H,T) &= \frac{1}{Z}\sum_{t \,:\, \mathrm{supp}(t) \subseteq T} g_H(t)\, S_t
  &&\text{(does one team hold the half-suit?)} \label{eq:q3}
\end{align}
with $g_H^{(c \to p)}$ the table restricted to survivors that give card $c$ to
seat $p$ and $t(A)$ the count vector of allocation $A$. All three are
\emph{exact} under the uniform prior of \S\ref{sec:constraints}; no
independence assumption and no sampling enters.

Equation~\eqref{eq:q2} is the one that matters for play. Declaring requires
naming an exact allocation, so the quantity that governs the decision is the
joint probability that all six assignments are simultaneously correct --- not
a product of per-card marginals, which is badly biased because the six cards
of a half-suit are strongly dependent through the capacity constraint
\eqref{eq:capacity}. \vthree{} scored declaration candidates with a geometric
mean of per-card marginals; \vfast{} and \vblock{} both target \eqref{eq:q2},
exactly in the latter case and approximately in the former (\S\ref{sec:fast}).

\begin{corollary}[Equal probability within a count vector]
\label{cor:equalprob}
Under the uniform posterior of \S\ref{sec:constraints}, any two allocations of
the same half-suit that satisfy (C1)--(C5) and share a count vector have equal
probability.
\end{corollary}

\begin{proof}
By \eqref{eq:q2} the probability of a surviving allocation $A$ depends on $A$
only through $t(A)$, since $S_t$ is a function of the count vector and the
dynamic-programming tables outside the block.
\end{proof}

The maximum-probability allocation is therefore determined by the count vector
alone, so the choice among survivors within a count vector is arbitrary and
\texttt{bestTeamAllocation} may return any of them. An allocation that
violates a certificate of \eqref{eq:disj}, however, has probability zero even
when its count vector is shared by survivors, so survival must be checked
explicitly rather than inferred from the count vector. The corollary and the
surviving-allocation caveat are both validated against exhaustive enumeration
on small reachable states in E15 (\S\ref{sec:results-verify}).

\subsection{A fast approximate posterior}
\label{sec:fast}

The reference engine is affordable per query, but six observers invoke it
after every public event, putting it in the inner loop of every large
experiment (\S\ref{sec:throughput}). We therefore also implement an
approximation, the \emph{Fast path}: it keeps the exact bookkeeping of
(C1)--(C3) and the exact capacities (C4), fits them by Sinkhorn scaling of the
support matrix, and interleaves a conditioning step \eqref{eq:cond} for (C5)
that reweights a certificate's cards $D$ by treating their entries $p_{d,a}$
as independent Bernoulli indicators before capacity fitting is re-run
(derivation and iteration schedule in \S\ref{app:sinkhorn}). The cost is
$O(\text{iters} \cdot |U| \cdot 6)$ rather than the $O(\prod_p (q_p+1))$ of
\eqref{eq:bound}.

Equation~\eqref{eq:cond} is an \emph{approximation}, and so are the $\hat\mu$,
$\hat\tau$ and $\hat\alpha$ it produces: its independence assumption is false,
because the capacity constraints that make \eqref{eq:q2} necessary also couple
the cards inside a certificate. E2 measures the gap over \numSelftestGames{}
games and \numSelftestChecks{} paired comparisons --- mean absolute difference
\numSinkhornMean{} between $\hat\mu$ and $\mu$, maximum \numSinkhornMax{},
accurate on average with a heavy tail whose positions E2 does not record
(\S\ref{sec:results-verify}). The Fast path is the compiled default
(\texttt{BeliefMode::Fast}) and produced every performance number reported in
this paper; that configuration is \vfast{}, and \vblock{} denotes the same
fitted policy with the exact reference engine substituted for it, which
appears only in the validation of \S\ref{sec:results-verify} and the
substitution ablation of \S\ref{sec:ablations}.

\subsection{The residual: policy-aware inference}
\label{sec:policyprior}

The posterior of (C1)--(C5) is exact with respect to the rules but is not the
full Bayesian posterior over deals, which would additionally weight each deal
$x$ by the likelihood $L(\mathcal{I} \mid x) = \prod_t \pi_{p_t}(a_t \mid
\mathcal{I}_{<t}, x)$ that the observed actions would have been chosen given
$x$. That likelihood is playstyle-dependent: a player who has taken many turns
and never asked in half-suit $H$ is, under any plausible policy, less likely
to hold cards of $H$, and the rules alone say nothing about this.

We implement a two-parameter family of such corrections as prior weights inside
the same machinery,
\begin{equation}
  w(c,p) \;=\; \exp\!\bigl(\theta\, a_{p,H(c)} - \phi\,(a_p - a_{p,H(c)})\bigr),
  \label{eq:prior}
\end{equation}
where $a_{p,H}$ counts $p$'s public asks in half-suit $H$, $a_p$ their total
public asks, and $H(c)$ is the half-suit of card $c$. The weights replace the
uniform initialisation of the Sinkhorn fit, so $\theta = \phi = 0$ recovers
the policy-agnostic posterior of \S\ref{sec:constraints}, and $\phi = 0$
recovers, in form, the heuristic \vthree{} used in place of the hard
certificate \eqref{eq:disj}. Equation \eqref{eq:prior} is \emph{fitted}:
$\theta$ and $\phi$ are two coordinates of the \numParams-coordinate vector
optimised jointly with the rest of the policy (\S\ref{sec:fitting}).

In the frozen-policy ablation E5 (\S\ref{sec:ablations}), disabling the prior ---
setting $\theta = \phi = 0$ while leaving every other coordinate at its fitted
value --- lowered the win rate by \numAblPrior{} points, 95\% paired percentile
bootstrap interval over deals \numAblPriorCI{}. The fitted policy therefore
depends on the prior. Because the remaining coordinates were fitted with the
prior active, this measures the sensitivity of \emph{this} fitted vector to
removing it and is not evidence about a policy retrained without it.
